\section{AMC 8 1986}
        \begin{enumerate}
        	\item (\textit{AMC 8 1986}) In July 1861, $366$ inches of rain fell in Cherrapunji, India.  What was the average rainfall in inches per hour during that month?


 $\text{(A)}\ \frac{366}{31\times 24} \qquad \text{(B)}\ \frac{366\times 31}{24}\qquad \text{(C)}\ \frac{366\times 24}{31}\qquad \text{(D)}\ \frac{31\times 24}{366}\qquad \text{(E)}\  366\times 31\times 24$

 \textbf{Answer: }A



\textbf{Solution 1}

July has $31$ days, or $31 \times 24$ hours.  Therefore, the average rainfall was $\boxed{\textbf{(A)}\ \frac{366}{31 \times 24}}$ inches per hour.

 



\textbf{Solution 2 (Dimensional Analysis)}

We can express the number of inches of rain in a \href{https://artofproblemsolving.com/wiki/index.php?title=Ratio}{ratio}, $\frac{366 \text{ inches}}{1\text{  month}}$. We need to find to find the number of hours, but first we need to find the number of days. We know there are 31 days in July, so we get 
 \[\frac{366 \text{ inches}}{1 \text{ month}} \times \frac{1 \text{ month}}{31 \text{ days}}\]

 We know there are 24 hours in a day, so we get 
 \[\frac{366 \text{ inches}}{1 \text{ month}} \times \frac{1 \text{ month}}{31 \text{ days}} \times \frac{1 \text{ day}}{24 \text{ hours}}\]

 And simplifying gets us $\dfrac{366 \text{ inches}}{31 \times 24 \text{ hours}}$, which is $\boxed{\text{A}}$.

 Note that in the simplifying step, we noted that the month and day units cancel, leaving us with only inches per hour.

 


	\item (\textit{AMC 8 1986}) Which of the following numbers has the largest reciprocal?


 $\text{(A)}\ \frac{1}{3} \qquad \text{(B)}\ \frac{2}{5} \qquad \text{(C)}\ 1 \qquad \text{(D)}\ 5 \qquad \text{(E)}\ 1986$

 \textbf{Answer: }A



\textbf{Solution 1}

Smaller positive real numbers have larger reciprocals.  Of these numbers, $\boxed{\textbf{(A)}\ \frac{1}{3} }$ is the smallest.

 



\textbf{Solution 2}

The reciprocals of the answer choices are $3$, $2\frac{1}{2}$, $1$, $\frac{1}{5}$, and $\frac{1}{1986}$.  Since $3$ is the largest, $\boxed{\textbf{(A)}\ \frac{1}{3} }$ is correct.

 


	\item (\textit{AMC 8 1986}) The smallest sum one could get by adding three different numbers from the set $\{ 7,25,-1,12,-3 \}$ is


 $\text{(A)}\ -3 \qquad \text{(B)}\ -1 \qquad \text{(C)}\ 3 \qquad \text{(D)}\ 5 \qquad \text{(E)}\ 21$

 \textbf{Answer: }C



\textbf{Solution}

The smallest possible sum is found by adding the $3$ smallest elements of the set.  This sum is  $(-3)+ (-1) + 7 = \boxed{\textbf{(C)}\ 3}$.

 


	\item (\textit{AMC 8 1986}) The product $(1.8)(40.3+.07)$ is closest to 


 $\text{(A)}\ 7 \qquad \text{(B)}\ 42 \qquad \text{(C)}\ 74 \qquad \text{(D)}\ 84 \qquad \text{(E)}\ 737$

 \textbf{Answer: }C



\textbf{Solution 1}

$(1.8)(40.3 + .07) = (1.8)(40.37)$, which is slightly more than $(1.8)(40) = 72$. So this expression is closest to $\boxed{\textbf{(C)}\ 74}$.

 -abcde26, edited by j314andrews

 



\textbf{Solution 2 (Brute Force)}

$(1.8)(40.3 + .07) = (1.8)(40.37) = 72.666$, which is closest to $\boxed{\textbf{(C)}\ 74}$.

 


	\item (\textit{AMC 8 1986}) A contest began at noon one day and ended $1000$ minutes later.  At what time did the contest end?


 $\text{(A)}\ \text{10:00 p.m.} \qquad \text{(B)}\ \text{midnight} \qquad \text{(C)}\ \text{2:30 a.m.} \qquad \text{(D)}\ \text{4:40 a.m.} \qquad \text{(E)}\ \text{6:40 a.m.}$

 \textbf{Answer: }D



\textbf{Solution}

The contest lasts $1000$ minutes, which is $\frac{1000}{60} = 16\frac{2}{3}$ hours or $16$ hours and $40$ minutes.  Since the contest starts at $\text{12:00 p.m.}$, it ends at $\boxed{\textbf{(D)}\  \text{4:40 a.m.}}$




\textbf{Solution 2}

There are $60$ minutes in an hour. So, we can easily eliminate some of the choices by noting that noon is exactly $720$ minutes away from midnight. Since $720 < 1000$, we know that it cannot be A or B. Because midnight is $720$ minutes away, we know that the contest ended $1000 - 720 = 280$ minutes after midnight. The highest multiple of 60 that will fit into $280$ is $240$, which is $4 \times 60$, and the remainder is $40$ minutes, meaning that the contest ended at $4:40 \text{ a.m.}$
$4:40$ is $\boxed{\text{D}}$



	\item (\textit{AMC 8 1986}) $\frac{2}{1-\frac{2}{3}}=$
$\text{(A)}\ -3 \qquad \text{(B)}\ -\frac{4}{3} \qquad \text{(C)}\ \frac{2}{3} \qquad \text{(D)}\ 2 \qquad \text{(E)}\ 6$

 \textbf{Answer: }E



\textbf{Solution}

$\frac{2}{1-\frac{2}{3}} = 2 \div \left(1 - \frac{2}{3}\right) = 2 \div \frac{1}{3} = 2 \cdot 3 = \boxed{\textbf{(E)}\ 6}$



	\item (\textit{AMC 8 1986}) How many whole numbers are between $\sqrt{8}$ and $\sqrt{80}$?


 $\text{(A)}\ 5 \qquad \text{(B)}\ 6 \qquad \text{(C)}\ 7 \qquad \text{(D)}\ 8 \qquad \text{(E)}\ 9$

 \textbf{Answer: }B



\textbf{Solution 1}

Note that $2 < \sqrt{8} < 3$ and $8 < \sqrt{80} < 9$.  Therefore, only the integers from $3$ to $8$ inclusive are between $\sqrt{8}$ and $\sqrt{80}$, for a total of $\boxed{\textbf{(B)}\ 6}$ integers.

 \textbf{Linear Approximation using Calculus} While this may not be the best, most efficient way to solve this problem, it can prove to be quite useful in similar problems. Using Calculus, we can use tangent-line approximation to approximate these square roots. First, we understand that these roots are of the form $f(x) = \sqrt{x}.$ Given the values to approximate, 8, and 80, we find the closest square number to them, namely, 9, and 81.

 Now, we take the derivative of $f(x)$ to get $f'(x) = \dfrac{1}{2\sqrt{x}}.$ We can approximate $f(x)$, to get $f(x) \approxeq f(x) + f`(x) \cdot \Delta{x},$ where $x$ is the perfect square closest to the the input value (in this case 8), and $\Delta{x}$ is the difference between $x$ and the input value.

 The two cases are as follows:

 
\begin{array}{|c|c|c|}input & x\text{ }value & \Delta{x} \\\hline8 & 9 & 8-9=-1 \\\hline80 & 81 & 80-81=-1 \\\end{array}

 Now, we can plug in our values into the equation, getting:

 
\begin{align*}&f(8) \approxeq f(9) + f`(9) \cdot -1 \implies \sqrt{9} + \dfrac{1}{2\sqrt{9}} \cdot -1  \implies 3 - \dfrac{1}{6} \\&f(80) \approxeq f(81) + f`(81) \cdot -1 \implies\sqrt{81} + \dfrac{1}{2\sqrt{81}} \cdot -1  \implies 9 - \dfrac{1}{18} \\\end{align*}


 Thus, we see that our answer lies in the amount of whole numbers between $3 - \dfrac{1}{6},$ and $9 - \dfrac{1}{18},$ and as we are counting whole numbers, we can round each expression to narrow our interval to $[3,8].$ Therefore, we see that our answer is $\boxed{6\text{ }(B)}.$

 Note: this solution is definitely not the fastest solution, and in a competition, you should probably use intuition to see the closest numbers in the interval of those square roots, and solve the problem from there.~Cheetahboy93

 


	\item (\textit{AMC 8 1986}) In the product shown, $\text{B}$ is a digit.  The value of $\text{B}$ is


 
 \[\begin{array}{rr} &\text{B}2 \\ \times &7\text{B} \\ \hline &6396 \\ \end{array}\]
$\text{(A)}\ 3 \qquad \text{(B)}\ 5 \qquad \text{(C)}\ 6 \qquad \text{(D)}\ 7 \qquad \text{(E)}\ 8$

 \textbf{Answer: }E



\textbf{Solution 1}

Recall that the units digit of the product of any two positive integers can be find by taking the units digit of each integer, multiplying these two digits, and taking the units digit of this product.

 In this case, $2 \times B$ has a units digit of $6$, so $B$ is either $3$ or $8$. If $B=3$, then the product is $32\times 73 < 40 \times 80 = 3200$, which is too small.  Thus $B = \boxed{\textbf{(E)}\  8}$, which works.

 



\textbf{Solution 2 (Quadratic Equation)}

This equation is equivalent to $(10B + 2)(70 + B) = 6396$.  Collecting all terms and factoring yields $10B^2 + 702B - 6256 = (B-8)(10B + 782) = 0$.  $B = \boxed{\textbf{(E)}\  8}$ is the only positive solution to this equation.

 


	\item (\textit{AMC 8 1986}) Using only the paths and the directions shown, how many different routes are there from $\text{M}$ to $\text{N}$?


 
                    \begin{figure}[H]
                        \centering
                        \includegraphics[width=0.5\textwidth]{Problem/AMC_8_AJHSME_1986_9_Problem_Figure_1.png}
                    \end{figure}
                    
$\text{(A)}\ 2 \qquad \text{(B)}\ 3 \qquad \text{(C)}\ 4 \qquad \text{(D)}\ 5 \qquad \text{(E)}\ 6$

 \textbf{Answer: }E



\textbf{Solution 1}

After one step, the possible routes are $M \rightarrow A$ and $M \rightarrow B$.

 After two steps, the possible routes are $M \rightarrow A \rightarrow D$, $M \rightarrow A \rightarrow C$, $M \rightarrow B \rightarrow A$, $M \rightarrow B \rightarrow C$, and $M \rightarrow B \rightarrow N$ (complete).

 After three steps, the possible routes are $M \rightarrow A \rightarrow D \rightarrow C$, $M \rightarrow A \rightarrow C \rightarrow N$ (complete), $M \rightarrow B \rightarrow A \rightarrow D$, $M \rightarrow B \rightarrow A \rightarrow C$, and $M \rightarrow B \rightarrow C \rightarrow N$ (complete).

 After four steps, the possible routes are $M \rightarrow A \rightarrow D \rightarrow C \rightarrow N$ (complete), $M \rightarrow B \rightarrow A \rightarrow D \rightarrow C$, and $M \rightarrow B \rightarrow A \rightarrow C \rightarrow N$ (complete).

 After five steps, the only possible route is $M \rightarrow B \rightarrow A \rightarrow D \rightarrow C \rightarrow N$ (complete).

 Therefore, there $\boxed{\textbf{(E)}\ 6}$ paths from $M$ to $N$.

 



\textbf{Solution 2}

There is $1$ path from $C$ to $N$, since the only edge leaving $C$ leads to $N$.  

 Since the only edge leaving $D$ leads to $C$, there is only $1$ path from $D$ to $N$.

 There are $2$ edges which leave $A$; one leads to $C$ and the other leads to $D$.  Therefore, there are $1 + 1 = 2$ paths from $A$ to $N$.

 There are $3$ edges which leave $B$; one leads to $A$, one leads to $C$, and one leads to $N$.  Therefore, there are $2 + 1 + 1 = 4$ paths from $B$ to $N$.

 There are $2$ edges which leave $M$; one leads to $A$ and the other leads to $B$.  Therefore, there are $4 + 2 = \boxed{\textbf{(E)}\ 6}$ paths from $M$ to $N$.

 The following diagram shows the number of paths to $N$ from each starting point.   

 
                    \begin{figure}[H]
                        \centering
                        \includegraphics[width=0.5\textwidth]{Solution/AMC_8_AJHSME_1986_9_Solution_Figure_1.png}
                    \end{figure}
                    

 \textbf{Note} This is an example of a directed acyclic graph (DAG); it has no cycles (paths that start and end at the same point).  A topological sort of this graph is $M, B, A, D, C, N$.

 


	\item (\textit{AMC 8 1986}) A picture $3$ feet across is hung in the center of a wall that is $19$ feet wide.  How many feet from the end of the wall is the nearest edge of the picture?


 $\text{(A)}\ 1\frac{1}{2} \qquad \text{(B)}\ 8 \qquad \text{(C)}\ 9\frac{1}{2} \qquad \text{(D)}\ 16 \qquad \text{(E)}\ 22$

 \textbf{Answer: }B



\textbf{Solution 1}

The wall is $19$ feet wide and the picture is $3$ feet wide, the uncovered parts of the wall total $16$ feet in width.  Since there is the same amount of empty space on each side of the picture, there is $16 \div 2 = \boxed{\textbf{(B)}\ 8}$ feet of space on each side of the picture.

 



\textbf{Solution 2}

Let $x$ be the distance in feet between an edge of the wall and the nearest edge of the picture. Since there are $x$ feet on both sides of picture, the total width of the wall is $2x + 3$ feet.  So $2x + 3 = 19$ and therefore $x = \boxed{\textbf{(B)}\ 8}$.

 


	\item (\textit{AMC 8 1986}) If $\text{A}*\text{B}$ means $\frac{\text{A}+\text{B}}{2}$, then $(3*5)*8$ is


 $\text{(A)}\ 6 \qquad \text{(B)}\ 8 \qquad \text{(C)}\ 12 \qquad \text{(D)}\ 16\qquad \text{(E)}\ 30$

 \textbf{Answer: }A



\textbf{Solution}

Evaluating yields $(3*5)*8 = \left( \frac{3+5}{2}\right) * 8 = 4*8 = \frac{4+8}{2} = \boxed{\textbf{(A)}\ 6}$.

 


	\item (\textit{AMC 8 1986}) The table below displays the grade distribution of the $30$ students in a mathematics class on the last two tests.  For example, exactly one student received a 'D' on Test 1 and a 'C' on Test 2 (see circled entry).  What percent of the students received the same grade on both tests?


 
                    \begin{figure}[H]
                        \centering
                        \includegraphics[width=0.5\textwidth]{Problem/AMC_8_AJHSME_1986_12_Problem_Figure_1.png}
                    \end{figure}
                    
$\text{(A)}\ 12\% \qquad \text{(B)}\ 25\% \qquad \text{(C)}\ 33\frac{1}{3}\% \qquad \text{(D)}\ 40\% \qquad \text{(E)}\ 50\%$

 \textbf{Answer: }D



\textbf{Solution}

Of the $30$ students, $2 + 4 + 5 + 1 + 0 = 12$ got the same grade on both tests.  So $\frac{12}{30} = 0.4 = \boxed{\textbf{(D)}\ 40\%}$ of the class got the same grade on both tests.

 


	\item (\textit{AMC 8 1986}) The perimeter of the polygon shown is 


 

                    \begin{figure}[H]
                        \centering
                        \includegraphics[width=0.5\textwidth]{Problem/AMC_8_AJHSME_1986_13_Problem_Figure_1.png}
                    \end{figure}
                    
$\text{(A)}\ 14 \qquad \text{(B)}\ 20 \qquad \text{(C)}\ 28 \qquad \text{(D)}\ 48$
$\text{(E)}\ \text{cannot be determined from the information given}$

 \textbf{Answer: }C



\textbf{Solution 1}

The unlabeled vertical sides must have total length $6$ while the unlabeled horizontal sides must have total length $8$.  Thus the unlabeled sides have total length $14$, so the total perimeter is $14 + 6 + 8 = \boxed{\textbf{(C)}\ 28}$.

 



\textbf{Solution 2}


                    \begin{figure}[H]
                        \centering
                        \includegraphics[width=0.5\textwidth]{Solution/AMC_8_AJHSME_1986_13_Solution_Figure_1.png}
                    \end{figure}
                    

 Extend two unlabeled sides of the original figure form a $6\times 8$ rectangle, as shown.  Notice that the extensions added to these two sides are the same lengths as the other two unlabeled sides of the original figure, so the $6\times 8$ rectangle has the same perimeter as the original figure.  Therefore, the perimeter of the original figure is $2(6+8) = \boxed{\textbf{(C)}\ 28}$.

 


	\item (\textit{AMC 8 1986}) If $200\leq a \leq 400$ and $600\leq b\leq 1200$, then the largest value of the quotient $\frac{b}{a}$ is


 $\text{(A)}\ \frac{3}{2} \qquad \text{(B)}\ 3 \qquad \text{(C)}\ 6 \qquad \text{(D)}\ 300 \qquad \text{(E)}\ 600$

 \textbf{Answer: }C



\textbf{Solution}

Since $a$ and $b$ are both positive, $\frac{b}{a}$ is maximized if $b$ is maximized and $a$ is minimized.  The maximum possible value of $b$ is $1200$ and the minimum possible value of $a$ is $200$.  In this case, $\frac{b}{a} = \frac{1200}{200} = \boxed{\textbf{(C)}\ 6}$.

 


	\item (\textit{AMC 8 1986}) Sale prices at the Ajax Outlet Store are $50\%$ below original prices.  On Saturdays an additional discount of $20\%$ off the sale price is given.  What is the Saturday price of a coat whose original price is $$180$?


 $\text{(A)}\ \$54 \qquad \text{(B)}\ \$72 \qquad \text{(C)}\ \$90 \qquad \text{(D)}\ \$108 \qquad \text{(E)}\ \$110$

 \textbf{Answer: }B



\textbf{Solution 1}

The sale price of the coat is $(1-0.5)(\$180) = (0.5)(\$180) = \$90$.  The Saturday price is $(1-0.2)(\$90) = (0.8)(\$90) = \boxed{\textbf{(B)}\ \$72}$.

 



\textbf{Solution 2}

After applying both discounts to an item, its price would be $(1-0.5)(1-0.2) = 0.5 \cdot 0.8 = 0.4$ times its original price.  Therefore, the Saturday price of the coat is $0.4 \cdot \$180 = \boxed{\textbf{(B)}\ \$72}$.

 


	\item (\textit{AMC 8 1986}) A bar graph shows the number of hamburgers sold by a fast food chain each season.  However, the bar indicating the number sold during the winter is covered by a smudge.  If exactly $25\%$ of the chain's hamburgers are sold in the fall, how many million hamburgers are sold in the winter?


 
                    \begin{figure}[H]
                        \centering
                        \includegraphics[width=0.5\textwidth]{Problem/AMC_8_AJHSME_1986_16_Problem_Figure_1.png}
                    \end{figure}
                    
$\text{(A)}\ 2.5 \qquad \text{(B)}\ 3 \qquad \text{(C)}\ 3.5 \qquad \text{(D)}\ 4 \qquad \text{(E)}\ 4.5$

 \textbf{Answer: }A



\textbf{Solution 1}

During the fall, $4$ million hamburgers were sold.  Since this is $25\%$, or $\frac{1}{4}$ of the company's hamburgers sold, a total of $4 \div \frac{1}{4} = 16$ million hamburgers were sold in total.  Since $4.5 + 5 + 4 = 13.5$ million hamburgers were sold in the other three seasons, $16 - 13.5 = \boxed{\textbf{(A)}\ 2.5}$ million hamburgers were sold in the winter.

 



\textbf{Solution 2}

Let $x$ be the number of hamburgers sold in the winter, in millions.Then a total of $4.5 + 5 + 4 + x = 13.5 + x$ million hamburgers were sold. Since $25\%$ of the hamburgers were sold in the fall, $0.25(13.5 + x) = 4$.  Solving this equation yields $x = \boxed{\textbf{(A)}\ 2.5}$.

 


	\item (\textit{AMC 8 1986}) Let $\text{o}$ be an odd whole number and let $\text{n}$ be any whole number.  Which of the following statements about the whole number $(\text{o}^2+\text{no})$ is always true?


 $\text{(A)}\ \text{it is always odd} \qquad \text{(B)}\ \text{it is always even}$
$\text{(C)}\ \text{it is even only if n is even} \qquad \text{(D)}\ \text{it is odd only if n is odd}$
$\text{(E)}\ \text{it is odd only if n is even}$

 \textbf{Answer: }E



\textbf{Solution 1}

Since $o$ is odd, $o^2$ must be odd.

 If $n$ is odd, then $no$ is odd, so $o^2 + no$ is even.  

 If $n$ is even, then $no$ is even, so $o^2 + no$ is odd.

 Therefore, of the answer choices, only $\boxed{\textbf{(E)}\ \text{it is odd only if } n \text{ is even}}$ accurately describes $o^2 + no$.

 



\textbf{Solution 2}

We are given that $\text{o}\equiv 1\pmod{2}$, so in mod $2$ we have 
 \[1^2+1(n) = n+1\] which is odd only if $\text{n}$ is even $\rightarrow \boxed{\text{E}}$




\textbf{Solution 3 (Answer Choices)}

If $o = 1$ and $n = 2$, then $o^2 + no = 3$, which eliminates (B) and (D).  If $o = 1$ and $n = 1$, then $o^2 + no = 2$, which eliminates (A) and (C).  This leaves $\boxed{\textbf{(E)}\ \text{it is odd only if } n \text{ is even}}$ as the correct answer.

 


	\item (\textit{AMC 8 1986}) A rectangular grazing area is to be fenced off on three sides using part of a $100$ meter rock wall as the fourth side.  Fence posts are to be placed every $12$ meters along the fence including the two posts where the fence meets the rock wall.  What is the fewest number of posts required to fence an area $36$ m by $60$ m?


 
                    \begin{figure}[H]
                        \centering
                        \includegraphics[width=0.5\textwidth]{Problem/AMC_8_AJHSME_1986_18_Problem_Figure_1.png}
                    \end{figure}
                    
$\text{(A)}\ 11 \qquad \text{(B)}\ 12 \qquad \text{(C)}\ 13 \qquad \text{(D)}\ 14 \qquad \text{(E)}\ 16$

 \textbf{Answer: }B



\textbf{Solution 1}

To minimize the total length of the fence, one of the longer ($60$ m) sides of the rectangle must touch the wall.  In this case, the total length of the fence is $36 + 60 + 36 = 132$ m.  Since pairs of consecutive posts are $12$ m apart, there are $132 \div 12 = 11$ intervals between consecutive posts, and therefore $11 + 1 = \boxed{\textbf{(B)}\ 12}$ posts are needed.  

 



\textbf{Solution 2}

As in Solution 1, one of the $60$ m sides must touch the wall.  The $60$ m side not touching the wall has $60 \div 12 + 1 = 6$ posts.  Each of the $36$ m sides has $36 \div 12 + 1 = 4$ posts.  Since $2$ posts are at intersections of $36$ m and $60$ m sides, there are a total of $6 + 2 \cdot 4 - 2 = \boxed{\textbf{(B)}\ 12}$ posts.

 


	\item (\textit{AMC 8 1986}) At the beginning of a trip, the mileage odometer read $56,200$ miles.  The driver filled the gas tank with $6$ gallons of gasoline.  During the trip, the driver filled his tank again with $12$ gallons of gasoline when the odometer read $56,560$.  At the end of the trip, the driver filled his tank again with $20$ gallons of gasoline.  The odometer read $57,060$.  To the nearest tenth, what was the car's average miles-per-gallon for the entire trip?


 $\text{(A)}\ 22.5 \qquad \text{(B)}\ 22.6 \qquad \text{(C)}\ 24.0 \qquad \text{(D)}\ 26.9 \qquad \text{(E)}\ 27.5$

 \textbf{Answer: }D



\textbf{Solution}

The car used $12$ gallons of fuel during the first part of the trip and $20$ gallons during the second part of the trip, for a total of $32$ gallons.  The total distance traveled during the trip was the difference between the final and initial odometer readings, that is, $57060 - 56200 = 860$ miles. Therefore, the car's fuel efficiency was $860 \div 32 \approx \boxed{\textbf{(D)}\ 26.9}$ miles per gallon.

 


	\item (\textit{AMC 8 1986}) The value of the expression $\frac{(304)^5}{(29.7)(399)^4}$ is closest to


 $\text{(A)}\ .003 \qquad \text{(B)}\ .03 \qquad \text{(C)}\ .3 \qquad \text{(D)}\ 3 \qquad \text{(E)}\ 30$

 \textbf{Answer: }D



\textbf{Solution}

By rounding, $\frac{(304)^5}{(29.7)(399)^4} \approx \frac{300^5}{30\cdot400^4} = \frac{3^5 \cdot 10^{10}}{3\cdot 4^4 \cdot 10^9} = \frac{3^4\cdot 10}{4^4} = \frac{810}{256} \approx \boxed{\textbf{(D)}\ 3}$.

 Note: The actual value of this expression, rounded to $8$ decimal places, is $3.44921198$.

 


	\item (\textit{AMC 8 1986}) Suppose one of the eight lettered identical squares is included with the four squares in the T-shaped figure outlined.  How many of the resulting figures can be folded into a topless cubical box?


 
                    \begin{figure}[H]
                        \centering
                        \includegraphics[width=0.5\textwidth]{Problem/AMC_8_AJHSME_1986_21_Problem_Figure_1.png}
                    \end{figure}
                    
$\text{(A)}\ 2 \qquad \text{(B)}\ 3 \qquad \text{(C)}\ 4 \qquad \text{(D)}\ 5 \qquad \text{(E)}\ 6$

 \textbf{Answer: }E



\textbf{Solution 1}

Suppose the existing net is folded so that the square immediately to the right of $D$ is on the bottom.  Then only the top and left sides of the folded cube would be missing.  $B$, $D$ and $F$ could be added as the left side of the cube, while $A$, $E$, and $H$ could be added as the right side of the cube.  $C$ and $G$ could not be added, as at most three faces can share a vertex.  Therefore, $\boxed{\textbf{(E)}\ 6}$ of the lettered faces could be added to the existing net.

 



\textbf{Solution 2}

Fold the four squares into the four sides of a cube. Then, there are six edges "open" (for lack of better term). For each open edge, we can add a square/side, so the answer is $6\rightarrow\boxed{\text{E}}$.

 


	\item (\textit{AMC 8 1986}) Alan, Beth, Carlos, and Diana were discussing their possible grades in mathematics class this grading period.  Alan said, "If I get an A, then Beth will get an A."  Beth said, "If I get an A, then Carlos will get an A."  Carlos said, "If I get an A, then Diana will get an A."  All of these statements were true, but only two of the students received an A.  Which two received A's?


 $\text{(A)}\ \text{Alan, Beth} \qquad \text{(B)}\ \text{Beth, Carlos} \qquad \text{(C)}\ \text{Carlos, Diana}$
$\text{(D)}\ \text{Alan, Diana} \qquad \text{(E)}\ \text{Beth, Diana}$

 \textbf{Answer: }C



\textbf{Solution}

If Alan got an A, then Beth, and therefore Carlos and Diana would also get A's.  But only two students actually got A's, so Alan did not get an A.

 If Beth got an A, then Carlos, and therefore Diana would get also get A's.  Once again, this is impossible since only two students got A's.  Therefore, Beth did not get an A.

 By process of elimination, $\boxed{\textbf{(C)}\ \text{Carlos, Diana}}$ must have gotten A's.

 


	\item (\textit{AMC 8 1986}) The large circle has diameter $\text{AC}$.  The two small circles have their centers on $\text{AC}$ and just touch at $\text{O}$, the center of the large circle.  If each small circle has radius $1$, what is the value of the ratio of the area of the shaded region to the area of one of the small circles?


 
                    \begin{figure}[H]
                        \centering
                        \includegraphics[width=0.5\textwidth]{Problem/AMC_8_AJHSME_1986_23_Problem_Figure_1.png}
                    \end{figure}
                    
$\text{(A)}\ \text{between }\frac{1}{2}\text{ and 1} \qquad \text{(B)}\ 1 \qquad \text{(C)}\ \text{between 1 and }\frac{3}{2}$
$\text{(D)}\ \text{between }\frac{3}{2}\text{ and 2} \qquad \text{(E)}\ \text{cannot be determined from the information given}$

 \textbf{Answer: }B



\textbf{Solution 1}

The small circle has radius $1$, thus its area is $\pi$.

 The large circle has radius $2$, thus its area is $4\pi$. 

 The area of the semicircle above $AC$ is then $2\pi$.

 The part that is not shaded are two small semicircles. Together, these form one small circle, hence their total area is $\pi$. This means that the area of the shaded part is $2\pi-\pi=\pi$. This is equal to the area of a small circle, hence the correct answer is $\boxed{\textbf{(B)}\ 1}$.

 



\textbf{Solution 2}

Let $A$ be the area of one of the small circles.  Since all circles are similar, the large circle's area is $2^2 \cdot A = 4A$.  So, the area inside the large circle but outside the small circles is $4A - 2A = 2A$.  Therefore, the shaded region must then be $\frac{2A}{2} = A$, so the desired ratio is $\frac{A}{A} = \boxed{\textbf{(B)}\ 1}$.

 


	\item (\textit{AMC 8 1986}) The $600$ students at King Middle School are divided into three groups of equal size for lunch.  Each group has lunch at a different time.  A computer randomly assigns each student to one of three lunch groups.  The probability that three friends, Al, Bob, and Carol, will be assigned to the same lunch group is approximately


 $\text{(A)}\ \frac{1}{27} \qquad \text{(B)}\ \frac{1}{9} \qquad \text{(C)}\ \frac{1}{8} \qquad \text{(D)}\ \frac{1}{6} \qquad \text{(E)}\ \frac{1}{3}$

 \textbf{Answer: }B



\textbf{Solution 1}

Suppose Al, Bob and Carol are the first three students to be randomly assigned groups, in that order.  After Al is assigned to a group, his group has $199$ open seats, out of a total of $599$ open seats across all the groups.  Therefore, the probability Bob is also assigned to this group is $\frac{199}{599}$, which is slightly less than $\frac{1}{3}$.  Similarly, if Al and Bob are assigned to the same group, the probability Carol is also assigned to this group is $\frac{198}{598}$, which is also slightly less than $\frac{1}{3}$.  So, the probability all three of them end up in the same group is $\frac{199}{599} \cdot \frac{198}{598}$, which is slightly less than $\frac{1}{3} \cdot \frac{1}{3} = \boxed{\textbf{(B)} \frac{1}{9}}$.

 


	\item (\textit{AMC 8 1986}) Which of the following sets of whole numbers has the largest average?


 $\text{(A)}\ \text{multiples of 2 between 1 and 101} \qquad \text{(B)}\ \text{multiples of 3 between 1 and 101}$
$\text{(C)}\ \text{multiples of 4 between 1 and 101} \qquad \text{(D)}\ \text{multiples of 5 between 1 and 101}$
$\text{(E)}\ \text{multiples of 6 between 1 and 101}$

 \textbf{Answer: }D



\textbf{Solution 1}

Let $f(n)$ be the average of the multiples of $n$ between $1$ and $101$.  There must be $\left\lfloor \frac{101}{n} \right\rfloor$ (the greatest integer less than or equal to $\frac{101}{n}$) multiples of $n$ between $1$ and $101$.  Let $g(n) = \left\lfloor \frac{101}{n} \right\rfloor$.  Then the multiples of $n$ between $1$ and $101$ must be $\{n, 2n, ..., g(n) \cdot n\}$, so their average is $f(n)=\frac{n + 2n + ... + g(n) \cdot n}{g(n)} = \frac{n}{g(n)} (1 + 2 + ... + g(n)) = \frac{n}{g(n)} \cdot \frac{g(n)(g(n)+1)}{2} = \frac{n(g(n)+1)}{2}$.  So $f(2) = 51$, $f(3) = 51$, $f(4) = 52$, $f(5) = 52.5$, and $f(6) = 51$, thus the set of $\boxed{\textbf{(D)}\ \text{multiples of 5 between 1 and 101}}$ has the greatest average.

 



\textbf{Solution 2}

The elements of each answer choice form arithmetic sequence. Since the sum of any $n$-term arithmetic sequence $\{a_1, ..., a_n\}$ is $\frac{n}{2}(a_1+a_n)$, its average is $\frac{1}{2}(a_1+a_n)$, that is, the average of its first and last terms.  Therefore, the averages of (A), (B), (C), (D), and (E) are $\frac{1}{2}(2+100) = 51$, $\frac{1}{2}(3+99) = 51$, $\frac{1}{2}(4+100) = 52$, $\frac{1}{2}(5+100) = 52.5$, and $\frac{1}{2}(6+96) = 51$, respectively. So of these, the set of $\boxed{\textbf{(D)}\ \text{multiples of 5 between 1 and 101}}$ has the largest average.

 


\end{enumerate}